\begin{abstract}
Legal retrieval-augmented generation (RAG) is usually evaluated as a single
accuracy number per benchmark, which obscures how it fails. We argue that
legal RAG failures are bottleneck-typed: across four legal benchmarks
(BarExam, HousingQA, CaseHOLD, LegalBench-SCALR) the useful intervention
changes with the dataset, the model, and the slice. We introduce a
\framework{} that labels each benchmark from cheap calibration traces (accuracy,
gold exposure, conditional accuracy, route disagreement, parse and error
counts, and average calls) and routes each row to the cheapest matching
intervention among baseline retrieval, legal query rewrite, \hyre{} (a
two-call HyDE variant in which the model commits to a tentative answer and
generates the controlling-rule passage that should support it before
retrieval), jurisdiction-scoped metadata filtering, option grounding,
conservative verifier policies, disagreement arbitration, and reject/escalate.
On Gemma~4~26B over four legal benchmarks, the controller lifts macro accuracy
from 71.9\% (matched baseline) to 77.9\% in calibration at 1.30 average LLM
calls per question, and from 71.5\% to 77.5\% on a held-out slice at 1.54
calls/Q. The lift is uneven and bottleneck-explained: a state-filter plus
conservative verifier route adds $+14\pp$ on HousingQA and a diverse-HyRE
route adds $+10\pp$ on CaseHOLD; BarExam exposes route-instability between
rewrite and \hyre{}, and SCALR's exact-match controller route ties baseline.
Fixed \hyre{} underperforms baseline on HousingQA (51.5\% vs.\ 60.5\%) and
CaseHOLD (72.0\% vs.\ 73.0\%); we therefore present it as one route inside a
portfolio whose value depends on diagnosis, not as a universal recipe. Our
contribution is the bottleneck taxonomy, the diagnostic protocol that types
each benchmark from its calibration trace, and the routing controller that
places generated reasoning on top of the active bottleneck.
\end{abstract}
